\section{Biological Example Family II: Working Memory and Attractor Maintenance}

One of the clearest reasons to import attractor language into this manuscript
comes from standard working-memory discussions in theoretical neuroscience.
Persistent activity and attractor-like maintenance models are useful because
they explain how a system can retain direction long enough to sustain a task
after the triggering cue has disappeared.

\begin{definition}[Integrator-like persistence]
An \textbf{integrator-like} mechanism is any structure that preserves task
relevant information long enough to bias subsequent motion in state space.
\end{definition}

\begin{discussion}
This matters for the book's framework in a practical way. A plan, note, or
external reminder can be read as a crude integrator-like aid. It preserves
directional information that would otherwise decay before the transition is
completed. The mathematics does not prove that a written note is literally a
cortical integrator, but it gives the reader a more serious concept than
``remember harder.''
\end{discussion}

\begin{example}[Writing session maintenance]
A graduate student opens a draft and writes three bullet claims before losing
momentum. If those bullet claims remain visible, they can serve as an
integrator-like memory trace that keeps the system near the writing basin.
Without them, the state may drift back toward reading or browsing. The working
memory analogy is therefore not a behavioural proof; it is a medium-lane
mathematical clarification of why visible persistence aids matter.
\end{example}
